\chapter*{Conclusion générale}
\addcontentsline{toc}{chapter}{Conclusion générale}
\markboth{\textbf{Conclusion générale}}{}

Ce projet de fin d'études au sein du Groupe COFAT a représenté une opportunité exceptionnelle de mettre en pratique et d'approfondir mes compétences en ingénierie logicielle dans un contexte industriel mondialisé et exigeant. La mission qui m'a été confiée — remplacer une gestion documentaire dispersée par un système d'information centralisé et temps réel — constituait un défi technique et organisationnel stimulant.

Grâce à une analyse rigoureuse de l'existant et à une collaboration étroite avec les départements méthodes et achats, j'ai pu concevoir, développer et déployer la plateforme "NVCapacity". Aujourd'hui, l'application remplit pleinement son rôle : elle permet aux responsables des usines (Tunisie, Mexique, Brésil) d'importer et de gérer numériquement leurs capacités locales en termes d'équipements, d'espaces industriels et de ressources humaines. Du côté de la direction et des achats, la plateforme garantit une consolidation instantanée et fiable, supprimant ainsi les lenteurs et les erreurs liées à l'agrégation manuelle de fichiers Excel. Par ailleurs, la dématérialisation du processus de budget non-industriel à travers un workflow de validation collaboratif a considérablement fluidifié les prises de décision.

Le succès de ce développement repose en grande partie sur l'adoption de la méthodologie Agile Scrum, qui m'a permis d'organiser mon travail en quatre sprints itératifs, d'assurer une livraison continue de valeur et de réajuster rapidement les fonctionnalités en fonction des retours du Product Owner. Sur le plan technique, l'utilisation d'une architecture 3-tiers moderne (Node.js, Express, React, SQL Server) m'a permis d'assurer la sécurité, la performance et l'évolutivité requises pour une application d'entreprise critique.

Ce stage de fin d'études a été particulièrement formateur. Il m'a non seulement permis de consolider mes compétences en développement Full-Stack, mais aussi d'acquérir une véritable compréhension des enjeux logistiques et financiers du secteur du câblage automobile.

Bien que la plateforme NVCapacity soit pleinement opérationnelle, de nombreuses perspectives d'amélioration peuvent être envisagées pour en étendre la portée. À court terme, il serait judicieux d'interfacer directement cette application avec l'ERP global du groupe (ex: SAP) afin d'automatiser la mise à jour de l'inventaire des machines physiques sans aucune saisie. À plus long terme, l'intégration de modèles d'Intelligence Artificielle et de Machine Learning sur l'historique des données de NVCapacity permettrait de faire de l'analyse prédictive, en anticipant les goulots d'étranglement capacitaires avant même qu'ils ne surviennent, optimisant ainsi encore davantage la stratégie d'investissement du Groupe COFAT.
